\documentclass[11pt]{article}
% ===========================================================================
%  Shared preamble for the AP Calculus AB exam-preparation worksheets.
%  Input right after \documentclass. Files beginning with "_" are not compiled
%  on their own. Navy and gold, to match the course's lesson pages.
% ===========================================================================
\usepackage[letterpaper,top=2.3cm,bottom=2.2cm,left=1.7cm,right=1.7cm,headheight=28pt,footskip=30pt]{geometry}
\usepackage{amsmath,amssymb}
\usepackage[T1]{fontenc}
\usepackage[utf8]{inputenc}
\usepackage{lmodern}
\usepackage{xcolor}
\usepackage[most]{tcolorbox}
\usepackage{enumitem}
\usepackage{fancyhdr}
\usepackage{array}
\usepackage{tikz}
\usepackage{pgfplots}
\pgfplotsset{compat=1.18}

\definecolor{navy}{HTML}{161E54}
\definecolor{ink}{HTML}{27357E}
\definecolor{gold}{HTML}{B98A00}
\definecolor{goldbg}{HTML}{FFFDF3}
\definecolor{inkbg}{HTML}{F6F7FD}
\definecolor{rule}{HTML}{C9CFF0}
\definecolor{okgreen}{HTML}{1B7A44}
\definecolor{okbg}{HTML}{F1F8F3}

\pagestyle{fancy}
\fancyhf{}
\fancyhead[L]{\footnotesize Name: \underline{\hspace{4cm}}}
\fancyhead[C]{\textbf{\color{navy}Integration Academy}\\[-2pt]{\scriptsize Calculus AB \textperiodcentered{} AP\textsuperscript{\textregistered} Exam Preparation}}
\fancyfoot[L]{\tiny\color{black!55}\parbox[t]{0.78\textwidth}{AP\textsuperscript{\textregistered} is a registered trademark of the College Board, which was not involved in the production of, and does not endorse, this material. \textcopyright{} Integration Academy.}}
\fancyfoot[R]{\scriptsize\thepage}
\renewcommand{\headrulewidth}{1.2pt}
\renewcommand{\headrule}{\hbox to\headwidth{\color{gold}\leaders\hrule height \headrulewidth\hfill}}
\renewcommand{\footrulewidth}{0pt}
\setlength{\parindent}{0pt}
\setlength{\parskip}{4pt}

% Title block: \wstitle{code}{title}{unit line}
\newcommand{\wstitle}[3]{\begin{center}{\small\bfseries\color{gold}WORKSHEET #1}\\[3pt]{\LARGE\bfseries\color{navy}#2}\\[3pt]{\small\color{black!60}#3}\end{center}\vspace{2pt}}
% Section heading
\newcommand{\sech}[1]{\par\vspace{10pt}{\large\bfseries\color{navy}#1}\par\nobreak\vspace{2pt}{\color{rule}\hrule height0.8pt}\vspace{7pt}}
% A level tag printed beside a group of questions
\newcommand{\level}[1]{\par\vspace{6pt}{\footnotesize\bfseries\color{gold}\MakeUppercase{#1}}\par\nobreak\vspace{2pt}}

\newtcolorbox{ideabox}{colframe=ink,colback=inkbg,coltitle=white,colbacktitle=ink,fonttitle=\bfseries\small,title={Key ideas},arc=2pt,boxrule=0.8pt,left=8pt,right=8pt}
\newtcolorbox{example}[1]{colframe=rule,colback=white,coltitle=navy,colbacktitle=inkbg,fonttitle=\bfseries\small,title={#1},arc=2pt,boxrule=0.8pt,borderline west={3pt}{0pt}{ink},left=10pt,right=8pt}
\newtcolorbox{frbox}[1]{colframe=gold,colback=goldbg,coltitle=white,colbacktitle=gold,fonttitle=\bfseries\small,title={#1},arc=2pt,boxrule=0.8pt,breakable,left=8pt,right=8pt}
\newtcolorbox{solbox}{colframe=okgreen,colback=okbg,coltitle=white,colbacktitle=okgreen,fonttitle=\bfseries\small,title={Worked solutions},arc=2pt,boxrule=0.8pt,breakable,left=8pt,right=8pt}
\newtcolorbox{scorebox}[1]{colframe=navy,colback=white,coltitle=white,colbacktitle=navy,fonttitle=\bfseries\small,title={#1},arc=2pt,boxrule=0.8pt,breakable,left=8pt,right=8pt}

\newcommand{\soln}{\par\smallskip\textbf{\color{ink}Solution.}\ }
\newcommand{\result}[1]{\par\smallskip\textbf{#1}}
% blank working space of a given height
\newcommand{\work}[1]{\par\vspace{#1}}
% a point value in the margin of a scoring guide
\newcommand{\pts}[1]{\hfill{\footnotesize\bfseries\color{gold}[#1]}}

% Coordinate grid for graph questions: \begin{apgraph}{xmin}{xmax}{ymin}{ymax} ... \end{apgraph}
\newenvironment{apgraph}[4]{\begin{center}\begin{tikzpicture}\begin{axis}[
  axis lines=middle, xmin=#1, xmax=#2, ymin=#3, ymax=#4,
  grid=both, grid style={black!12}, xtick distance=1, ytick distance=1,
  ticklabel style={font=\scriptsize}, xlabel={$x$}, ylabel={$y$},
  width=10cm, height=7.2cm, samples=80, clip=false,
  every axis plot/.append style={line width=1.1pt, color=ink}]}{\end{axis}\end{tikzpicture}\end{center}}
\newcommand{\closeddot}[2]{\addplot[only marks, mark=*, mark size=2.2pt, color=ink] coordinates {(#1,#2)};}
\newcommand{\opendot}[2]{\addplot[only marks, mark=*, mark size=2.2pt, color=ink, mark options={fill=white}] coordinates {(#1,#2)};}

% Multiple-choice options: four across, or stacked when the options are long.
\newcommand{\opts}[4]{\par\vspace{3pt}\begin{tabular}{@{}p{0.235\linewidth}p{0.235\linewidth}p{0.235\linewidth}p{0.235\linewidth}@{}}(A)\ #1 & (B)\ #2 & (C)\ #3 & (D)\ #4\end{tabular}\par\vspace{9pt}}
\newcommand{\optsv}[4]{\par\vspace{3pt}\begin{tabular}{@{}l@{\ \ }p{0.9\linewidth}@{}}(A) & #1\\(B) & #2\\(C) & #3\\(D) & #4\end{tabular}\par\vspace{9pt}}

\fancyhead[R]{\footnotesize\color{black!60}Worksheet 0.1}
\begin{document}

\wstitle{0.1}{Functions You Must Know Cold}{Unit 0: Readiness}

\begin{ideabox}
\begin{itemize}[leftmargin=*,itemsep=2pt,topsep=2pt]
  \item A \textbf{domain} excludes zero denominators, negative numbers under even roots, and non-positive inputs of logarithms.
  \item Changes \emph{inside} the brackets act on $x$ and work backwards: $f(x-h)$ moves the graph right $h$. Changes \emph{outside} act on $y$.
  \item $f\big(g(x)\big)$ means $g$ first, then $f$. Its domain is found before simplifying.
  \item To find an \textbf{inverse}, write $y=f(x)$, solve for $x$, and swap the letters. The domain of $f^{-1}$ is the range of $f$.
\end{itemize}
\end{ideabox}

\sech{Worked examples}
\begin{example}{Example 1 \textperiodcentered{} Two restrictions at once}
Find the domain of $f(x)=\dfrac{\sqrt{2x+8}}{x-3}$.\soln The root needs $2x+8\ge0$, that is $x\ge-4$. The denominator is zero at $x=3$.\result{Domain: $x\ge-4$ and $x\neq3$.}
\end{example}

\begin{example}{Example 2 \textperiodcentered{} Composition in both orders}
Let $f(x)=\dfrac1x$ and $g(x)=x-4$. Find $f\big(g(x)\big)$ and $g\big(f(x)\big)$ with their domains.\soln $f\big(g(x)\big)=\dfrac{1}{x-4}$, defined for $x\neq4$. $g\big(f(x)\big)=\dfrac1x-4$, defined for $x\neq0$.\result{The two compositions are different functions with different domains.}
\end{example}

\begin{example}{Example 3 \textperiodcentered{} Inverse of a rational function}
Find the inverse of $f(x)=\dfrac{x-1}{x+2}$.\soln $y(x+2)=x-1$ gives $xy-x=-2y-1$, so $x=\dfrac{-2y-1}{y-1}=\dfrac{2y+1}{1-y}$.\result{$f^{-1}(x)=\dfrac{2x+1}{1-x}$.}
\end{example}

\newpage
\sech{Practice}

\level{Core}
\begin{enumerate}[leftmargin=*,itemsep=2pt]
  \item Find the domain of $f(x)=\sqrt{10-5x}$.\work{2.3cm}
  \item State the range of $y=|x|-3$.\work{2.3cm}
  \item Let $f(x)=4-x^2$ and $g(x)=2x+1$. Find $f\big(g(1)\big)$ and $g\big(f(1)\big)$.\work{2.3cm}
  \item Find the inverse of $f(x)=5x-3$.\work{2.3cm}
\end{enumerate}

\level{Exam level}
\begin{enumerate}[leftmargin=*,itemsep=2pt,start=5]
  \item Find the domain of $f(x)=\ln\!\left(x^2-4\right)$.\work{2.8cm}
  \item Let $f(x)=\sqrt{x}$ and $g(x)=9-x^2$. Find $f\big(g(x)\big)$ and state its domain.\work{2.8cm}
  \item Find the inverse of $f(x)=e^{2x}+1$, and state the domain of the inverse.\work{2.8cm}
  \item Describe how the graph of $y=\dfrac{3}{x+2}-1$ is obtained from the graph of $y=\dfrac1x$, and write equations for its two asymptotes.\work{2.8cm}
\end{enumerate}

\begin{frbox}{Question 9 \textperiodcentered{} Free-response practice}
Let $f(x)=x^2-6x$ with domain $x\ge3$.
\begin{enumerate}[label=(\alph*),leftmargin=*,itemsep=2pt]
  \item Find the range of $f$.\work{2.6cm}
  \item Find $f^{-1}(x)$. (Completing the square is helpful.)\work{2.6cm}
  \item State the domain and range of $f^{-1}$, and verify that $f^{-1}\big(f(5)\big)=5$.\work{2.6cm}
\end{enumerate}
\end{frbox}

\newpage
\begin{solbox}
\begin{enumerate}[leftmargin=*,itemsep=4pt]
  \item $10-5x\ge0$, so $x\le2$.
  \item $y\ge-3$.
  \item $g(1)=3$ and $f(3)=-5$. $f(1)=3$ and $g(3)=7$.
  \item $f^{-1}(x)=\dfrac{x+3}{5}$.
  \item $x^2-4>0$, so $x<-2$ or $x>2$.
  \item $\sqrt{9-x^2}$, with domain $-3\le x\le3$.
  \item $y-1=e^{2x}$ gives $x=\dfrac12\ln(y-1)$, so $f^{-1}(x)=\dfrac12\ln(x-1)$ with domain $x>1$.
  \item Stretched vertically by 3, moved left 2 and down 1. Asymptotes: $x=-2$ and $y=-1$.
  \item \begin{enumerate}[label=(\alph*),leftmargin=*,itemsep=3pt]
    \item $f(x)=(x-3)^2-9$, which is least at $x=3$: range $y\ge-9$.
    \item $y+9=(x-3)^2$ with $x\ge3$ gives $x=3+\sqrt{y+9}$, so $f^{-1}(x)=3+\sqrt{x+9}$.
    \item Domain $x\ge-9$, range $y\ge3$. $f(5)=-5$ and $f^{-1}(-5)=3+\sqrt4=5$.
  \end{enumerate}
\end{enumerate}
\end{solbox}

\end{document}
